\documentclass[../main.tex]{subfiles}
\begin{document}
%%%%%%%%%%%%%%%%%%%%%%%%%%%%%%%%%%%%%%%%%%%%%%%%%%%%%%%%%%%%%%%%%%%%%%%%%%%%%%%%%%%%%%%%%%%%%%%%%%%
\chapter{Managing Canvases, Layers, and Objects in the Sidebar}

Most of the work you do in OmniGraffle is on the canvas. That's where you draw things, but
\textit{managing} your OmniGraffle project happens mostly in the Sidebar. Positioned to the
left of the canvas, the Sidebar serves many purposes.

The Sidebar is divided into four different tabs, with the \textbf{Layers} tab being the primary
view. You use the \textbf{Layers} tab to keep track of the canvases, layers, and objects in
your project, but there's lots more you can do here; for example:

\begin{itemize}
  \item You can use the \textbf{Guides} tab to add, position, and change the colors of ruler
    guides.
  \item You can use the \textbf{Outline} tab to enter structured, hierarchical outlines to
    create diagrams and flowcharts.
  \item You can use the \textbf{Selection} tab to view the objects in your project by a variety
    of styles, making it easier for you to mass-select objects and work on them independently
    from everything else on the canvas.
\end{itemize}

This chapter provides you with a thorough look at OmniGraffle's new Unified Sidebar.

\section{Meet the (New) Sidebar}

One of the big UI refinements you'll discover is the new \textbf{Sidebar}. The sidebar helps you
manage your work, providing you with quick access to the canvases, layers, and objects in your
project in one convenient location.

Along the top of the Sidebar, you'll see four tabs for getting at every aspect of your project.
By default, new projects expose the \textbf{Layers} tab in the Sidebar.

\ogfigure{0.800}{res/og7_sidebar01.png}

To switch between the different tabs, you can either click the buttons along the top of the
Sidebar, or you can use the following keyboard shortcuts:

\begin{itemize}
  \item Layers (\textbf{Option-Command-2})
  \item Guides (\textbf{Option-Command-3})
  \item Outline (\textbf{Option-Command-4})
  \item Selection (\textbf{Option-Command-5})
\end{itemize}

You have many options for hiding and showing the Sidebar:

\begin{itemize}
  \item Click the Hide/Show Sidebar button in the toolbar
  \item Choose \textbf{View} $\blacktriangleright$ \textbf{Hide/Show Sidebar} in the menu bar
  \item Use the \textbf{Option-Command--1} keyboard shortcut
\end{itemize}

\begin{ogtip}
  Similarly, you can hide and show the Inspector bar on the other side of the canvas by pressing
  \textbf{Shift-Command-I}. These two keyboard shortcuts come in handy when you want to hide
  everything while working on the canvas.
\end{ogtip}

You can also choose to have more of a compact view in the Sidebar, which is particularly helpful
when you have a project with a lot of canvases or layers. \textbf{Control-/Right-click} in the
Sidebar, and choose \textbf{Use small icons} from the contextual menu that appears:

\ogfigure{0.800}{res/og7_sidebar_smallicons.png}

\begin{ognote}
  The Sidebar's contextual menu also contains the handy \textbf{Show Canvas Page Count} option,
  which toggles display of the number of printed pages each canvas requires. This number is
  based on page dimensions set in Page Setup, Canvas Size and Units set in the Canvas Inspector,
  and Margins set in the Document Inspector.
\end{ognote}

Use the \textbf{Search} field at the bottom of the Sidebar to isolate objects by a specific
string of text. As you enter characters, OmniGraffle starts narrowing down the selection so
you can find what you need, when you need it. Click the Search field's magnifying glass to
isolate the search criteria to text found in the object's Name, Notes, Data, or Text, and to
examine any recent searches you have conducted.

\ogfigure{0.673}{res/og7_sidebar02.png}

The following sections provide you with an overview of the four Sidebar tabs.

\subsection{The Layers Tab \ogicon{1em}{res/og7_icon_layersview_selected.png}}

The Layers tab is likely where you will conduct most of your Sidebar-related work. The Layers
tab lists everything in your OmniGraffle project: canvases, layers, and everything you've
drawn on those layers.

\subsubsection{Differentiating Between Layer Types in the Sidebar}

Depending on which version of OmniGraffle you have---Standard or Pro---there are three different
types of layers that you can add to your project:

\ogfigure{0.673}{res/og7_sidebar_layers01.png}

\begin{itemize}
  \item \textbf{Standard Layer} --- The most basic form of Layer; every canvas you create has at
    least one standard layer. Standard layers are unique to each canvas in your project;
    meaning, their contents cannot be shared with other canvases.
\end{itemize}

\begin{ogpro}
  \begin{itemize}
    \item \textbf{Shared Layer} (\textbf{PRO}) --- Shared Layers are layers that you can share
      across all of the canvases in your project. Shared Layers, when empty, have an
      orangish-brown background color in the sidebar. Similarly, the objects on a Shared Layer
      take on the same orangish-brown color as a highlight. Use a Shared Layer in projects where
      you need a repeating element across the different canvases, such as adding a header or
      footer, or adding printers' crop marks to your design.
    \item \textbf{Artboard Layer} (\textbf{PRO}) --- Artboard Layers have many different uses
      (as you'll see in Using Artboards (PRO)), but the key thing to know is that they affect
      objects on layers \textit{above} the Artboard Layer. Artboard Layers, when empty, have a
      red background color in the sidebar. Similarly, the objects on an Artboard Layer take on
      the same red color as a highlight color.
  \end{itemize}
\end{ogpro}

\subsubsection{A Guide to the Layer Icons}

When you look at the Layers tab in the Sidebar, you'll notice a series of icons next to each
layer. These icons provide you with a quick visual reference as to that layer's current state.

Let's take a closer look.

\ogfigure{0.800}{res/og7_sidebar_layers02.png}

\begin{itemize}
  \item Immediately to the left of the Canvas and Layer name is a square \textit{preview} image.
    The \textbf{Canvas Preview} shows everything that's on all visible layers in that canvas;
    the \textbf{Layer Preview} shows the objects on that particular layer.
  \item Each layer has a three icons to the right of the layer name:

    \begin{itemize}
      \item \textbf{Visible or Hidden}: \ogicon{8.3pt}{res/og7_icon_layer_visible.png} means that
        the layer is visible; a layer is hidden (not visible) when its icon is
        \ogicon{12.5pt}{res/og7_icon_layer_hidden.png}. Objects on \textbf{Hidden} layers are not
        visible in the Canvas Preview, but they are, however, visible on the Layer Preview.
        Additionally, hidden layers are not exported or printed.
      \item \textbf{Printable or Not-Printable}: \ogicon{12.5pt}{res/og7_icon_layer_print.png}
        means that the objects on this layer are printable; or, not printable if you see this icon
        \ogicon{12.5pt}{res/og7_icon_layer_dontprint.png}. In the Export panel, you can opt to
        allow non-printing layers to be included as part of the export.
      \item \textbf{Unlocked or Locked}: \ogicon{12.5pt}{res/og7_icon_layer_unlocked.png}
        indicates that the layer is unlocked and objects on that layer can be moved around. When
        the layer is locked \ogicon{12.5pt}{res/og7_icon_layer_locked.png}, the objects on that
        layer cannot be moved and you cannot add anything else to that layer until the layer is
        unlocked.
    \end{itemize}
  \item When you are working on a layer, the layer will have a
    \ogicon{12.5pt}{res/og7_icon_layer_pencil_unlocked.png} to the left of its disclosure
    triangle. If the layer is locked and you select the layer, you'll see
    \ogicon{12.5pt}{res/og7_icon_layer_pencil_locked.png}, which indicates that you cannot add
    to or change anything on that layer until it is unlocked
    \ogicon{12.5pt}{res/og7_icon_layer_unlocked.png}.
  \item If you have an object selected on an editable layer, the pencil gets a little object box
    around it \ogicon{12.5pt}{res/og7_icon_layer_editable.png}.
\end{itemize}

\begin{ognote}
  For more information about Shared and Artboard Layers, see Sharing Layers with Multiple Canvases
  (PRO) or Using Artboards (PRO).
\end{ognote}

\subsubsection{Adding a New Layer to a Canvas}

There are many options for adding a new layer to a Canvas:

\begin{itemize}
  \item Click the \textbf{New Layer} button in the toolbar
  \item Choose \textbf{Edit} $\blacktriangleright$ \textbf{Layers} $\blacktriangleright$
    \textbf{New Layer} from the menu bar
  \item In the Layers sidebar, right-click on the Canvas and choose \textbf{New Layer} from the
    contextual menu that appears
  \item Right-click on the Canvas and choose \textbf{Layers} $\blacktriangleright$ \textbf{New
    Layer} from the contextual menu that appears
  \item Click \ogicon{12.5pt}{res/og7_button_actionmenu.png} at the bottom of the Sidebar and
    choose \textbf{New Layer} from the Action menu
\end{itemize}

\begin{ognote}
  Before adding a new layer to your project, you may wish to switch to the \textbf{Layers} tab in
  the Sidebar, by either choosing \textbf{View} $\blacktriangleright$ \textbf{Show Contents}
  $\blacktriangleright$ \textbf{Layers}, or by pressing \textbf{Option-Command-2}.
\end{ognote}

\subsubsection{Moving and Duplicating Layers}

To move a layer from one canvas to another, simply click-and-drag the layer from one canvas to
the next.

If what you want is a duplicate of the layer, do one of the following:

\begin{itemize}
  \item Select the layer in question and choose \textbf{Edit} $\blacktriangleright$
    \textbf{Layers} $\blacktriangleright$ \textbf{Duplicate Layer}
  \item Control-/Right-click on the layer and choose \textbf{Duplicate Layer} from the
    contextual menu that appears
  \item Control-/Right-click on the canvas and choose \textbf{Layers} $\blacktriangleright$
    \textbf{Duplicate Layer} from the contextual menu that appears
  \item Click \ogicon{12.5pt}{res/og7_button_actionmenu.png} at the bottom of the Sidebar and
    choose \textbf{Duplicate Layer} from the Action menu
\end{itemize}

\subsubsection{Merging Layers}

If your project has multiple layers and you'd like to combine the shapes and objects on two
layers onto a single layer, you can do that by selecting the layers and choosing \textbf{Edit}
$\blacktriangleright$ \textbf{Layers} $\blacktriangleright$ \textbf{Merge Layer Down}.

As the menu option implies, OmniGraffle relies on the layer stacking order when merging layers.
The layers must be on top of each other in the Layers sidebar; you cannot merge non-contiguous
layers.

\ogfigure{0.800}{res/og7_sidebar_layersXX.png}

At left in the previous image, the objects on Layer 3 would merge down into Layer 2. However, if
you were to choose \textbf{Merge Layer Down} with the selection shown on the right, Layer 5
would merge into Layer 4---and Layers 3 and 2 would remain unchanged.

\subsubsection{Deleting a Layer}

To delete a layer from a Canvas:

\begin{itemize}
  \item Select the layer in question and press the \textbf{Delete} key
  \item Select the layer in question and choose \textbf{Edit} $\blacktriangleright$
    \textbf{Layers} $\blacktriangleright$ \textbf{Delete Layer}
  \item Control-/Right-click on the layer and choose \textbf{Delete Layer} from the contextual
    menu that appears
  \item Control-/Right-click on the canvas and choose \textbf{Layers} $\blacktriangleright$
    \textbf{Delete Layer} from the contextual menu that appears
  \item Click \ogicon{12.5pt}{res/og7_button_actionmenu.png} at the bottom of the Sidebar and
    choose \textbf{Delete Layer} from the Action menu
\end{itemize}

And remember, if you delete a layer by accident, you can always hit \textbf{Command-Z} to
restore the layer.

\begin{ognote}
  A Canvas must have at least one layer, so you can never delete every single layer in a project.
\end{ognote}

\subsubsection{Changing Canvas, Layer, and Object Names in the Sidebar}

As you add new objects, layers, and canvases to your projects, you'll notice that their names in
the sidebar aren't all that informative. For example, new projects start with \textbf{Canvas
1} as the canvas name, or \textbf{Layer 1} as that canvas's first layer. Create a new layer
for that canvas, and it's name is \textbf{Layer 2} and so on. Add a new canvas to the project,
and it gets named \textbf{Canvas 2}.

\ogfigure{0.800}{res/og7_sidebar_naming01.png}

These names are just placeholders, and they await your thing-labeling skills to assign them
names that make sense for the project in which you are working. For example, to change the
name of \textbf{Canvas 1} to something else, double-click on the Canvas Name to make the name
editable, and then enter something logical, such as \textbf{Project Workflow} or \textbf{My
Diabolical Robot}.

\ogfigure{0.800}{res/og7_sidebar_naming02.png}

The same theory applies to Layer Names; you'll want to change the names of individual layers so
that each one denotes its purpose in the project. For example, to change the name of
\textbf{Layer 1}, double-click on the Layer Name to make the name editable, and then enter
something logical, such as \textbf{Basic Shapes} or \textbf{Left Foot}.

\ogfigure{0.800}{res/og7_sidebar_naming03.png}

And using a similar process, you can change the name of the shapes or objects on the layers of
your project. No longer are you stuck with a project full of \textbf{Rectangle}s; you can
change those to \textbf{Squares} or \textbf{Robot's Big Toe}, if you'd like.

\ogfigure{0.800}{res/og7_sidebar_naming04.png}

If you enter a long name for something, as shown in the previous image, that's okay. OmniGraffle
truncates the name for display purposes only; however, the width of the sidebar is expandable.
Click in the space between the sidebar and the canvas's vertical ruler, and then drag left or
right as needed.

\ogfigure{0.800}{res/og7_sidebar_naming05.png}

Naming the canvases, layers, and objects helps you to quickly identify where things are in your
project.

\subsection{The Guides Tab \ogicon{1em}{res/og7_icon_guidesview_selected.png}}

OmniGraffle comes with a full complement of guides that you can use to perfectly align objects
on the canvas. There are Ruler Guides, Smart Alignment Guides, and Smart Distance Guides;
however, only the Ruler Guides---the guides which you can drag onto the canvas from the left
and top ruler bars---can be managed from the Sidebar.

\begin{ognote}
  To learn more about Smart Alignment Guides and Smart Distance Guides, see Basic Concepts,
  earlier in this guide.
\end{ognote}

You can use a combination of Ruler Guides and OmniGraffle's Smart Guides to position and
perfectly align objects on the canvas. To add a Ruler Guide to your canvas, click and hold on
either the top or left ruler bar and then drag a guide onto the canvas.

\ogfigure{0.673}{res/og7_sidebar_guides01.png}

You can also click the Action menu at the bottom of the Sidebar, and choose \textbf{New
Horizontal Guide} or \textbf{New Vertical Guide}, depending on which one you need.

\ogfigure{0.800}{res/og7_sidebar_guides02.png}

The Ruler Guides are bright pink by default, but you can change the color by clicking the pop-up
menu to the right of each guide in the Sidebar. When you do, the Colors window appears,
allowing you to change the color to anything you'd like.

\ogfigure{0.800}{res/og7_sidebar_guides03.png}

There's no limit to the number of Ruler Guides you can drag onto the canvas, so don't be shy.

\begin{ogtip}
  You can also change the guide colors using the \textbf{Appearance} preference pane. For more
  details, see Appearance, earlier in this guide.
\end{ogtip}

\subsubsection{Hiding Ruler Guides}

Okay, so you took our advice and you weren't shy about using Ruler Guides. You've used plenty of
them\ldots{}and that's okay.

But now you need to see what your masterpiece looks like without all of the guides in the way.
Easy-peasy! To hide the Ruler Guides, click one of the rulers, and to make them reappear,
click the rulers again.

\subsubsection{Duplicating Guides from an Existing Canvas to a New Canvas}

When you add a new canvas to a project, that canvas adopts the Ruler Guides of the previously
selected canvas. For example, if you have Ruler Guides set up on Canvas 1 and you click
\textbf{New Canvas} in the toolbar, Canvas 2 appears with the same Ruler Guides applied.

\ogfigure{0.769}{res/og7_sidebar_guides05.png}

However, Ruler Guides can be canvas-independent; meaning, you can have different Ruler Guides on
each canvas.

\subsubsection{Copying Ruler Guides from One Canvas to Another}

Let's say that you've been working on a project with multiple canvases, each with multiple Ruler
Guides. You decide to copy a newly created shape to another canvas, and there's a Ruler Guide
that you'd like to go along with it. You could always drag out a new Ruler Guide on that
canvas, but then you need to move it into position. Instead, you can select a Ruler Guide that
exists on one canvas and drag it into the canvas that you'd like it on.

\ogfigure{0.769}{res/og7_sidebar_guides06.png}

\subsubsection{Deleting a Ruler Guide}

To delete a Ruler Guide, you can drag the guide back toward the ruler and it disappears. If you
are using an \textbf{Infinite Canvas}, be sure to drag the guide to the ruler, since there are
no bounds for the canvas itself. You can also delete Ruler Guides from the Guides tab by
clicking to select the guide you want to delete from the canvas, and then press the
\textbf{Delete} key.

\ogfigure{0.800}{res/og7_sidebar_guides07.png}

\subsection{The Outline Tab \ogicon{1em}{res/og7_icon_outlineview_selected.png}}

\ogfigure{0.350}{res/og7_sidebar_outline01.png}

The Outline tab gives you a quick way to enter structured data to create hierarchical diagrams.
The Outline shows any shapes you've created on the canvas, arranged hierarchically according
to the lines that connect them.

Likewise, as you enter text into an Outline, shapes are automatically created and connected on
the canvas with a corresponding text label. To help facilitate the flow of creating an
outline, however, you should turn on \textbf{Auto layout} in the \textbf{Diagram Layout and
Style} inspector.

\ogfigure{0.800}{res/og7_sidebar_outline02.png}

To create an entirely new outline, click the Action menu at the bottom of the Sidebar and choose
\textbf{New Topic}.

\ogfigure{0.800}{res/og7_sidebar_outline03.png}

While you work in the Outline tab, the following keys can be used to create and arrange items:

\begin{itemize}
  \item \textbf{Return} --- creates a ``peer'' of the currently selected item. Peers appear on
    the same indentation level in the outline, and, if they are beyond the top level, have
    connections from the same parent object in the diagram.
  \item \textbf{Tab} --- indents the currently selected item, making it a ``child'' of the item
    above it in the outline. A connection line is drawn in the diagram from the parent object to
    the child object.
  \item \textbf{Shift-Tab} --- outdents the currently selected item, making it a peer of what
    used to be its parent.
  \item \textbf{Escape (esc)} --- switches between selecting an item and editing the item's
    text.
\end{itemize}

\begin{ognote}
  Additional options and keyboard shortcuts for quickly building outlines can be found in
  \textbf{Edit} $\blacktriangleright$ \textbf{Outlining}.
\end{ognote}

You can rearrange items in the Outline tab by selecting and dragging the item to another
location in the Outline. Dragging an item with children brings its children along for the
ride. While you drag an item, a blue line shows the location and indentation level at which
the item will be dropped when you let go of the mouse button.

\ogfigure{0.800}{res/og7_sidebar_outline04.png}

A group, table, or subgraph on the canvas is shown as a single item in the Outline. To view the
contents of the group, table, or subgraph in the Outline, you will need to ungroup those
objects first.

While the Outline Editor is handy for quickly entering lots of information, an outline cannot
always perfectly represent the relationships in a complicated diagram. Objects with more than
one parent, lines connected to other lines, and circular ancestry are all examples of the
kinds of things a simple outline can't properly show. In these cases, the Outline tab tries
its best to be accurate. If your diagram is not strictly tree-like, you probably shouldn't
rely too much on the Outline if what you need is a quick representation of the objects'
relationships.

\subsection{The Selection Tab \ogicon{1em}{res/og7_icon_selectionview_selected.png}}

The Selection tab gives you the ability to cloak or reveal the different elements in your
project, making it easy for you to examine and apply styles to the objects in your project.

\ogfigure{0.800}{res/og7_sidebar_selection01.png}

Along the top of the Selection tab, you will see a row of tiny \textbf{Style Selection} buttons
that look oddly familiar. Each button corresponds to the styles you can set for the shapes on
the Canvas.

\ogfigure{0.800}{res/og7_sidebar_selection02.png}

All of these are on by default, and this gives you the finest level of granularity when
evaluating your objects. But as you turn one---or all---of them off, you'll notice that the
previews in the Selection tab become more compact.

\ogfigure{0.800}{res/og7_sidebar_selection03.png}

Beneath that row of tiny buttons, you'll also find a pop-up menu, which lets you select
individual canvases in your project, or \textbf{All Canvases} for getting the full scope of
your project.

\ogfigure{0.800}{res/og7_sidebar_selection04.png}

\subsection{Sharing Layers with Multiple Canvases (\texorpdfstring{\textbf{PRO}}{PRO})}

\begin{ogpro}
  You can share a layer between multiple canvases. That way, whenever you change the layer, it
  affects every canvas that uses that layer.

  To create a new Shared Layer, do one of the following:

  \begin{itemize}
    \item Choose \textbf{Edit} $\blacktriangleright$ \textbf{Layers} $\blacktriangleright$
      \textbf{New Shared Layer}
    \item In the Layers sidebar, Control-/Right-click on the desired canvas and choose
      \textbf{New Shared Layer} from the contextual menu that appears
    \item Control-/Right-click on the canvas and choose \textbf{Layers} $\blacktriangleright$
      \textbf{New Shared Layer} from the contextual menu that appears
  \end{itemize}

  These methods create a Shared Layer on the current canvas only. To apply the new Shared Layer
  to other canvases, click to select it in the Layers sidebar and drag it to the other
  canvases where you want it to show up.

  To convert an existing layer to a Shared Layer:

  \begin{itemize}
    \item Choose \textbf{Edit} $\blacktriangleright$ \textbf{Layers} $\blacktriangleright$
      \textbf{Layer Settings} $\blacktriangleright$ \textbf{Shared Layer}
    \item In the Layers sidebar, Control-/Right-click on an existing layer and choose
      \textbf{Layer Settings > Shared Layer} from the contextual menu that appears
  \end{itemize}

  These methods convert the selected layer to a Shared Layer and apply it to all canvases in the
  project immediately.

  Shared Layers have an orange tint added to the Layer Preview, and the objects on the Shared
  Layer also gain an orange highlight color when selected on the canvas.

  After adding a new Shared Layer:

  \begin{itemize}
    \item Add and edit objects on the layer normally.
    \item Any further changes you make to any copy of the layer are also applied to all other
      copies of it.
  \end{itemize}

  To delete a Shared Layer, do one of the following:

  \begin{itemize}
    \item Select the Shared Layer and press the \textbf{Delete} key
    \item Choose \textbf{Edit} $\blacktriangleright$ \textbf{Delete} from the menu Bar
    \item Control-/Right-click on the Shared Layer and select \textbf{Delete Layer} in the
      contextual menu
    \item Click the \ogicon{12.5pt}{res/og7_button_actionmenu.png} and select \textbf{Delete
      Layer} from the Action menu at the bottom of the sidebar
  \end{itemize}

  Deleting a Shared Layer in this way only removes the currently selected instance of the layer.

  To remove all instances of the layer from the document, Control-/Right-click on one of them in
  the Layer sidebar and choose \textbf{Delete (layer name) from All Canvases}, or choose
  \textbf{Delete (layer name) from All Canvases} from the Action menu at the bottom of the
  sidebar.
\end{ogpro}
%%%%%%%%%%%%%%%%%%%%%%%%%%%%%%%%%%%%%%%%%%%%%%%%%%%%%%%%%%%%%%%%%%%%%%%%%%%%%%%%%%%%%%%%%%%%%%%%%%%
\end{document}
